% Liste des notations et abréviations, placée après les listes (structure du mémoire de référence).
\chapter*{Liste des notations et abréviations}
\addcontentsline{toc}{chapter}{Liste des notations et abréviations}
{\renewcommand{\arraystretch}{0.92}%
\begin{tabular}{@{}ll@{}}
AEVB & Bayes variationnel auto-encodant (\emph{auto-encoding variational Bayes}) \\
AIC & critère d'information d'Akaike \\
ARI & indice de Rand ajusté \\
BIC & critère d'information bayésien \\
CAVI & montée par coordonnées pour l'inférence variationnelle \\
DM & test de Diebold et Mariano \\
ELBO & borne inférieure de l'évidence (\emph{evidence lower bound}) \\
EM & espérance-maximisation (\emph{expectation-maximization}) \\
EM-BIC & EM dont le nombre de composantes est choisi par le BIC \\
EM-ICL & EM dont le nombre de composantes est choisi par l'ICL \\
FEWS NET & réseau des systèmes d'alerte précoce contre la famine \\
GEM & EM généralisé \\
GMM & modèle de mélange gaussien \\
HMM & modèle de Markov caché \\
IC & intervalle de confiance \\
ICL & critère de vraisemblance complétée intégrée \\
KL & divergence de Kullback-Leibler \\
LDA & allocation de Dirichlet latente \\
MAP & maximum a posteriori \\
MAR & données manquantes aléatoires (\emph{missing at random}) \\
MCAR & données manquantes complètement aléatoires (\emph{missing completely at random}) \\
MCMC & méthodes de Monte-Carlo par chaînes de Markov \\
MLE & estimateur du maximum de vraisemblance \\
NLP & traitement automatique du langage naturel \\
NPMI & information mutuelle ponctuelle normalisée \\
PAM & Programme alimentaire mondial \\
pLSA & analyse sémantique latente probabiliste \\
PPCA & analyse en composantes principales probabiliste \\
SEM & EM stochastique \\
SGVB & estimateur bayésien variationnel par gradient stochastique \\
SVI & inférence variationnelle stochastique \\
TOST & procédure des deux tests unilatéraux (\emph{two one-sided tests}) \\
VAE & auto-encodeur variationnel \\
VB & Bayes variationnel (EM variationnel en champ moyen) \\
VI & inférence variationnelle \\
\end{tabular}}
